%% ============================================================
%% 第一章 引言（中文版）
%% 结构：P1 研究背景+问题引出 → P2 用 Quantum Kernel 解决上述问题
%%       → P3 Our Work → P4 主要贡献（3 条）+ 本文结构
%% 编译：xelatex（ctex）
%% 说明：全文遵循"用英文逻辑写中文"，句句有主谓宾、句间逻辑显式。
%%       引用示例统一采用"xxxx年，某某等提出……，通过……，实现了……"的句式。
%% ============================================================

\section{引言}

时间序列预测是一类依据历史观测推断系统未来演化状态的建模任务，其核心目标是从数据中学习系统的动态规律，为后续的调度与控制提供提前判断的依据。
作为深度学习中的一个重要方向，它把深度神经网络与时序分析相融合，形成了一种新型的智能预测范式，显著提升了预测精度。
近年来，时间序列预测已被广泛用于电力调度、金融市场与交通管理等领域。
实际系统通常由多个变量共同描述，变量之间彼此关联，多变量时间序列预测因此成为主要的研究场景，其关键在于刻画变量间的相互依赖关系。
面向多变量预测，许多研究者针对时序建模提出了不同的网络架构与改进方法，预测精度不断提升。
例如，2024年，Liu 等提出 iTransformer，把注意力机制作用于变量维，以每个变量作为 token 建模变量间关系，在多变量预测上取得了显著效果\cite{liu2024itransformer}；2025年，Chittoor 等提出 QuLTSF，把经典线性网络与变分量子线路相结合，首次把量子机器学习用于长时序预测\cite{chittoor2025qultsf}；2024年，Gu 和 Dao 提出 Mamba，以选择性状态空间模型实现线性复杂度的序列建模，成为长序列建模的重要方法\cite{gu2024mamba}。
Mamba 是一类基于状态空间模型的序列建模架构，它通过输入相关的选择机制对状态进行动态更新，在保持线性复杂度的同时刻画长程依赖，因而避免了自注意力机制的二次复杂度。
在 Mamba 架构之上，最具代表性的多变量时间序列预测模型是 S-Mamba\cite{wang2025smamba}。S-Mamba 把多变量序列的每个变量视为一个 token，以变量维为序列维进行正向与反向两次状态空间扫描，再把两个方向的输出融合，以线性复杂度完成多变量预测。
针对 S-Mamba，一系列优化方法被相继提出。例如，2025年，Ma 等提出 TimePro，在变量扫描过程中自适应选择关键时间步调制隐状态，以更少的参数量取得更好的预测效果\cite{ma2025timepro}；2026年，Karadag 等提出 ms-Mamba，让多个采样率不同的状态空间模型并行工作，从多个时间尺度刻画序列\cite{karadag2026msmamba}；2025年，Jura 等提出 Q-SSM，用变分量子门的期望值调节状态空间模型的记忆更新，改善了 S-Mamba 训练不稳定、对初始化敏感的问题\cite{jura2025qssm}。
尽管针对 S-Mamba 的优化方法不断被提出，其跨变量耦合建模能力仍未被充分挖掘。
上述方法分别从关键时间步的选择、多时间尺度的建模与记忆更新的稳定性等方面对 S-Mamba 加以改进，但都未改变变量间耦合只能被隐式携带这一事实。
由于双向结构沿变量维逐位传递信息，变量之间的关系只能被隐式混入各变量的隐状态之中，模型既无法分离出任一变量对之间的依赖，也无法给出耦合的强度数值。
这使得模型无法区分强耦合变量与弱耦合变量，也难以在预测时抑制无关变量带来的干扰，对跨变量输入的响应因而极弱，长程预测精度随之受限。
更重要的是，由于耦合关系始终以隐状态的形式存在，模型依据哪些变量间关系作出预测无法被直接读取，预测结果缺乏可核对的依据。
因此，构建一种能够显式刻画变量间耦合、并让这份耦合直接参与预测的建模方法，已成为多变量时间序列预测中亟待解决的重要问题。

幸运的是，我们发现量子核方法可以为解决上述挑战提供思路。
量子核方法以量子态之间的重叠程度衡量两条数据的相似程度，它首先把数据经量子电路编码为量子态，再以两个量子态的内积作为核值，核值越大表示两条数据在量子特征空间中越相似。
由于量子态的振幅与相位都可以携带信息，量子核能够表达经典核函数难以刻画的高维非线性关系，因而具备建模数据之间耦合关系的潜力。
近些年，利用量子核计算耦合关系的机器学习工作相继被提出。
例如，2024年，Zhao 等提出量子核自注意力网络 QKSAN，以量子核替代自注意力中的相似度计算，使注意力权重在量子特征空间中得到表达\cite{zhao2024qksan}；2024年，Bai 等提出层次对齐的量子 Jensen-Shannon 核，以量子态之间的相似度度量图结构之间的对应关系，在图表分类任务上取得良好效果\cite{bai2024haqjsk}。
这些研究表明，量子核具备建模耦合关系的能力。
因此，本文提出一种基于量子核的时序编码方法，用于有效表征数据、反映变量之间的关系，从而解决 S-Mamba 在跨变量建模上的不足。

本文提出量子跨变量耦合核 QCCK，把多变量时序中变量间的耦合显式地构造出来。
QCCK 由频域对齐、量子编码与量子核三个环节构成。
频域对齐作为 QCCK 的输入处理环节，在时间轴与频率轴上分别消除变量间的时移与周期尺度差异，把每个变量的原始序列变换为对齐频谱。
量子编码通过振幅通道与相位通道把每个变量的信息编码为量子态，并借助参数化量子线路在比特间建立纠缠，得到反映变量特征的纠缠量子态。
量子核以两个量子态的重叠程度度量变量间的耦合强度，对任意两个变量输出一个取值在 0 到 1 之间的数值，全部变量对的结果构成耦合矩阵 $K$，其中每个元素都表示一对变量的耦合强度，可以直接读取。
在 QCCK 的基础上，本文进一步提出面向多变量时间序列预测的模型 QCCK-M。QCCK-M 把 QCCK 以耦合层的形式嵌入主干编码器层之间，使显式度量得到的耦合矩阵参与跨变量信息融合，从而在逐层的表示更新中利用变量之间的显式依赖关系完成预测。
实验结果表明，QCCK-M 在多个公开数据集上优于现有方法，其中在 Beijing 数据集上 MSE 降低 11\%，在 ETTm2 上降低 6.2\%，在 ECL 上降低 6.1\%。
同时，耦合矩阵以显式数值呈现变量间的耦合强度，模型所依据的变量关系可以被直接读取与核对。
QCCK 的意义在于把变量间的耦合从隐式传递变为显式可读，模型依据哪些变量间关系作出预测，可以直接检查与核对。
更重要的是，QCCK-M 让这份显式耦合真正参与逐层的变量表示更新，使可解释性不止停留在事后分析，而成为预测过程本身的依据。

本文的主要贡献如下：

01 提出量子跨变量耦合核 QCCK，它以频域对齐消除变量间的时移与周期尺度差异，以振幅通道与相位通道把变量编码为纠缠量子态，并以量子核度量变量间的耦合强度，从而显式地构造出可解释的耦合矩阵 $K$，解决了变量间耦合难以被显式刻画与利用的问题。

02 构建多变量时间序列预测网络 QCCK-M，它整合输入预处理、编码器层的逐层更新、跨变量耦合融合与预测输出，把耦合矩阵作为跨变量信息融合的依据，使显式耦合直接参与变量表示更新，形成完整的端到端预测模型。

03 在多组公开数据集上验证了 QCCK-M 的有效性，其中在 Beijing 数据集上 MSE 降低 11\%，在 ETTm2 上降低 6.2\%，在 ECL 上降低 6.1\%；并结合耦合矩阵的可解释性分析，说明模型所依据的变量关系可以被直接读取与核对。

本文余下部分组织如下：

第 2 节介绍相关研究与量子计算基础。
第 3 节提出量子跨变量耦合核 QCCK。
第 4 节提出多变量时间序列预测网络 QCCK-M。
第 5 节给出实验设置、主结果、消融与可解释性分析。
最后，第 6 节得出结论。
